\input{../tex-inputs/latex-leseansicht-vorspann}

\section[Theodor von Sosnosky an Arthur Schnitzler, 26. 5. 1901]{L04231 Theodor von Sosnosky an Arthur Schnitzler, 26. 5. 1901}
\nopagebreak\mylabel{L04231v}
\rehead{ }\normalsize\beginnumbering\briefempfaengerindex{Schnitzler, Arthur@\textsc{Schnitzler, Arthur}!zzzSosnosky, Theodor von@\emph{von Theodor von Sosnosky}!1901-05-262@{26. 5. 1901}|(be}
\toendnotes[C]{\smallbreak\pagebreak[2]}
\correspDesc{Versand  durch Theodor von Sosnosky am 26. 5. 1901 in Kremsmünster
\newline{}Erhalt  durch Arthur Schnitzler im Zeitraum [27. 5. 1901
                  – 31. 5. 1901?] in Wien}\toendnotes[C]{\smallbreak}
\Standort{DLA, A:Schnitzler, HS.1985.1.4640.}
\physDesc{Brief, 2 Blätter, 8 Seiten, 6614 Zeichen
\newline{}Handschrift: blaue Tinte, deutsche Kurrent
\newline{}Schnitzler: mit Bleistift datiert: »«901«  }\toendnotes[C]{\smallbreak}
\pstart
           {\pb}dzt \textsc{Kremsmünster Ob-Oesterr.\oindex{Kremsmünster@\textbf{Kremsmünster}, \emph{Verwaltungsgebiet}|pw}}\pend
           
\pstart
           26./5\pend
           
\pstart{}Sehr geehrte Herr Doktor!\pend\vspace{0.5em}
\pstart
           Ihr eben eingetroffener frdl. Brief veranlaßt mich, ihn gleich zu beantworten, weil
               ich ihm zu entnehmen glaube, daß meine Äußerung im Dankſchreiben nicht ganz richtig
               verſtanden worden iſt. Ich beeile mich daher, Ihnen hier über das, was ich dort nur
               flüchtig angedeutet, ausführlichere Aufklärung zu geben:\pend
           
\pstart
           Selbstverständlich fällt es mir gar nicht ein, einem Autor des Recht, sein Themen zu
               wählen, ſtreitig machen zu wollen; um ſo weniger, alſ ich der Anſicht bin, in der
               Kunſt müße {\pb}\uline{Alles} \strikeout{w} erlaubt ſein,
               wenn es nur künstleriſch dargeſtellt ſei; (eine Ausnahme mache ich blos für die
               Bühne, auf der ich \introOben{}ſozial\introOben{} aufreizende Stücke zu geben für
               äußerst thöricht halte, Stücke wie die Weber\pwindex{\textcolor{red}{\textsuperscript{XXXX indx1}}!Weber. Schauspiel aus den vierziger Jahren@\strich\emph{Die Weber. Schauspiel aus den vierziger Jahren}|pw}
               z. B. werden nur als Aufreizung \strikeout{vom} zum Klaßenhaß
               empfunden und \label{K_L04231-7v}\edtext{ausgeſchrotet}{\lemma{\textnormal{\emph{ausgeschrotet}}}\Cendnote{\textnormal{östrreichisch veraltet: ausgeschlachtet, zum Verkauf in Teile zerlegt, ausgenutzt}}}\label{K_L04231-7}).\pend
           
\pstart
           Die \uline{künſtlerischen} Qualitäten Ihrer Novelle\pwindex{Schnitzler, Arthur 15.\,5.\,1862 Wien – 21.\,10.\,1931 ebd.@\textsc{Schnitzler, Arthur} (15.\,5.\,1862 Wien – 21.\,10.\,1931 ebd.), \emph{Schriftsteller, Mediziner}!Lieutenant Gustl. Novelle@\strich\emph{Lieutenant Gustl. Novelle}|pwv} will und wollte ich alſo gar nicht
               antaſten, im Gegenteil: ich kann nur aus vollſter Überzeugung wiederholen, ich finde
               die Darſtellung techniſch und pſychologisch \uline{brillant}{ }{\kaufmannsund}{ }ſtelle das Werk\pwindex{Schnitzler, Arthur 15.\,5.\,1862 Wien – 21.\,10.\,1931 ebd.@\textsc{Schnitzler, Arthur} (15.\,5.\,1862 Wien – 21.\,10.\,1931 ebd.), \emph{Schriftsteller, Mediziner}!Lieutenant Gustl. Novelle@\strich\emph{Lieutenant Gustl. Novelle}|pwv} künſtleriſch \uline{höher} als
                  »Garlan\pwindex{Schnitzler, Arthur 15.\,5.\,1862 Wien – 21.\,10.\,1931 ebd.@\textsc{Schnitzler, Arthur} (15.\,5.\,1862 Wien – 21.\,10.\,1931 ebd.), \emph{Schriftsteller, Mediziner}!Frau Bertha Garlan. Roman@\strich\emph{Frau Bertha Garlan. Roman}|pw}«, ja als Alles, was \strikeout{Sie} ich biſher von Ihnen kenne, den köſtlichen »Reigen\pwindex{Schnitzler, Arthur 15.\,5.\,1862 Wien – 21.\,10.\,1931 ebd.@\textsc{Schnitzler, Arthur} (15.\,5.\,1862 Wien – 21.\,10.\,1931 ebd.), \emph{Schriftsteller, Mediziner}!Reigen. Zehn Dialoge@\strich\emph{Reigen. Zehn Dialoge}|pw}« – Ihr Meisterstück – ausgenommen.\pend
           
\pstart
           Anderseits hielt ich es aber für geboten, in meinem Briefe anzudeuten, daß ich
               perſönlich dieſer Arbeit\pwindex{Schnitzler, Arthur 15.\,5.\,1862 Wien – 21.\,10.\,1931 ebd.@\textsc{Schnitzler, Arthur} (15.\,5.\,1862 Wien – 21.\,10.\,1931 ebd.), \emph{Schriftsteller, Mediziner}!Lieutenant Gustl. Novelle@\strich\emph{Lieutenant Gustl. Novelle}|pwv}{ }{\pb}\label{K_L04231-1v}\edtext{\textsc{\begin{otherlanguage}{french}vis à vis\end{otherlanguage}}}{\lemma{\textnormal{\emph{vis à vis}}}\Cendnote{\textnormal{französisch: von Angesicht zu
                  Angesicht}}}\label{K_L04231-1} sozusagen oppoſitionell empfinde, weil ich darin eine – Ihnen
               vielleicht unbewußte {\kaufmannsund} von Ihnen ungewollte Gehäßigkeit
               gegen die Offiziere erblicke und dieser meiner perſönlichen Empfindung in einer
               eventuellen Beſprechung des Buches\pwindex{Schnitzler, Arthur 15.\,5.\,1862 Wien – 21.\,10.\,1931 ebd.@\textsc{Schnitzler, Arthur} (15.\,5.\,1862 Wien – 21.\,10.\,1931 ebd.), \emph{Schriftsteller, Mediziner}!Lieutenant Gustl. Novelle@\strich\emph{Lieutenant Gustl. Novelle}|pwv} auch Ausdruck geben würde. Diese Bedenken ſcheinen mir die
               künſtlerischen Vorzüge des Buches\pwindex{Schnitzler, Arthur 15.\,5.\,1862 Wien – 21.\,10.\,1931 ebd.@\textsc{Schnitzler, Arthur} (15.\,5.\,1862 Wien – 21.\,10.\,1931 ebd.), \emph{Schriftsteller, Mediziner}!Lieutenant Gustl. Novelle@\strich\emph{Lieutenant Gustl. Novelle}|pwv} doch nicht zu tangiren?!\pend
           
\pstart
           Aus zwei Bemerkungen glaube ich die Gehäßigkeit – wie gesagt: die unbewußte –
               deutlich herauszuleſen; ich werde dieſe gegebenen Falles in der Beſprechung erwähnen.
               Bevor ich dieſe ſchreibe, will ich das Buch\pwindex{Schnitzler, Arthur 15.\,5.\,1862 Wien – 21.\,10.\,1931 ebd.@\textsc{Schnitzler, Arthur} (15.\,5.\,1862 Wien – 21.\,10.\,1931 ebd.), \emph{Schriftsteller, Mediziner}!Lieutenant Gustl. Novelle@\strich\emph{Lieutenant Gustl. Novelle}|pwv} noch Andern zu leſen geben, deren Urteil mir wert iſt.
               Ein ſolches – vor meinem gefälltes – sti{\geminationm}t mit meinem
               ganz überein. Meine Empfindung ſcheint demnach keineswegs vereinzelt zu ſtehen. Die
                  \label{K_L04231-2v}\edtext{dem »\textsc{Echo\pwindex{literarische Echo@\emph{Das literarische Echo}|pw}}« entno{\geminationm}ene Nachricht}{\lemma{\textnormal{\emph{dem … Nachricht}}}\Cendnote{\textnormal{»Gegen Arthur
                        Schnitzler, der Oberarzt im nichtaktiven Stande der österreichischen\oindex{Österreich-Ungarn@\textbf{Österreich-Ungarn}|pw} Landwehr ist, wurde wegen
                     Veröffentlichung seiner Offiziersnovelle ›Leutnant Gustl\pwindex{Schnitzler, Arthur 15.\,5.\,1862 Wien – 21.\,10.\,1931 ebd.@\textsc{Schnitzler, Arthur} (15.\,5.\,1862 Wien – 21.\,10.\,1931 ebd.), \emph{Schriftsteller, Mediziner}!Lieutenant Gustl. Novelle@\strich\emph{Lieutenant Gustl. Novelle}|pw}‹ vom Landwehroberkommando Untersuchung
                     eingeleitet.«, \emph{Nachrichten. Allerlei}. In: \emph{Das litterarische Echo}\pwindex{literarische Echo@\emph{Das literarische Echo}|pwk}, Jg. 3, Nr. 9, Februar 1901, Sp. 650.}}}\label{K_L04231-2}, daß Ihnen dieſe Novelle\pwindex{Schnitzler, Arthur 15.\,5.\,1862 Wien – 21.\,10.\,1931 ebd.@\textsc{Schnitzler, Arthur} (15.\,5.\,1862 Wien – 21.\,10.\,1931 ebd.), \emph{Schriftsteller, Mediziner}!Lieutenant Gustl. Novelle@\strich\emph{Lieutenant Gustl. Novelle}|pwv} auch Unannehmlichkeiten in Ihrer militäriſchen Stellung \strikeout{durch} als Landwehr-Arzt {\pb}zugezogen hat,
               iſt abermals ein Argument, das für mich ſpricht. Ich will damit aber nicht etwa
               behaupten, daß ich das Vorgehen der milit. Behörden \uline{gutheiße}; aber \uline{begreifen} kann ich’s. Gewiß
               liegen in Ihrer Arbeit\pwindex{Schnitzler, Arthur 15.\,5.\,1862 Wien – 21.\,10.\,1931 ebd.@\textsc{Schnitzler, Arthur} (15.\,5.\,1862 Wien – 21.\,10.\,1931 ebd.), \emph{Schriftsteller, Mediziner}!Lieutenant Gustl. Novelle@\strich\emph{Lieutenant Gustl. Novelle}|pwv} keine
               thatſächlichen \label{K_L04231-3v}\edtext{\textsc{gravamina}}{\lemma{\textnormal{\emph{gravamina}}}\Cendnote{\textnormal{lateinisch: Beschwerden}}}\label{K_L04231-3}, aber –
                  »\label{K_L04231-4v}\edtext{\begin{otherlanguage}{french}\textsc{c’est le ton qui fait la musique}\end{otherlanguage}}{\lemma{\textnormal{\emph{c’est … musique}}}\Cendnote{\textnormal{französisch: der Ton macht die
                  Musik}}}\label{K_L04231-4}« und daß dieſer Ton gegenüber dem Militär kein freundlicher iſt,
               haben eben auch die Behörden empfunden {\kaufmannsund} daraus ihre
               Konſequenzen gezogen, die wieder einmal eklatant beweiſen, daß die Inſtitution der
               Reſerveoffiziere, wie ſie bei uns beſteht, unzulänglich iſt {\kaufmannsund} zu den schwerſten Kolliſionen führen kann, weil
                  die \introOben{}ſozialen\introOben{} Anſchauungen des Privatmannes denen des
               Offiziers oft ſtracks zuwiderlaufen; ein Fall, der in Deutschland\oindex{Deutschland@\textbf{Deutschland}|pw} unmöglich ist, weil dort \uline{nur} der
               Reſerve-Offizier wird, von dem man poſitiv weiß, daß er keinerlei Geſi{\geminationn}ung hegt, die mit der eines Offiziers kollidirt. –
                Verzeihen Sie gütigſt, daß ich dieſe Dinge erwähne, aber da {\pb}Ihrer
                  Novelle\pwindex{Schnitzler, Arthur 15.\,5.\,1862 Wien – 21.\,10.\,1931 ebd.@\textsc{Schnitzler, Arthur} (15.\,5.\,1862 Wien – 21.\,10.\,1931 ebd.), \emph{Schriftsteller, Mediziner}!Lieutenant Gustl. Novelle@\strich\emph{Lieutenant Gustl. Novelle}|pw} nun einmal auch außer der
               litterarisch-künſtlerischen eine \introOben{}andere\introOben{} Bedeut\strikeout{t}ung zugemeßen worden \strikeout{ist} und dies in die Öffentlichkeit gedrungen \introOben{}iſt\introOben{}, ſo konnte ich nicht umhin, bei der durch Ihren Brief veranlaßten
               Erörterung \strikeout{die} auch dieſe Sache zu beleuchten.
               Was ſpeziell Ihren Einwand betrifft, \substVorne{}\textsuperscript{falls}\substDazwischen{}wenn\substHinten{} in dieſem Falle Gehäſſigkeit gegen die Offiziere vorläge,{ }ſo \substVorne{}\textsuperscript{ſei}\substDazwischen{}wäre\substHinten{} dann \introOben{}auch\introOben{} ſolche gegen Liberale, Ärzte,
               Weinhändler, Theater-Direktoren im »\textsc{Vermächtnis\pwindex{Schnitzler, Arthur 15.\,5.\,1862 Wien – 21.\,10.\,1931 ebd.@\textsc{Schnitzler, Arthur} (15.\,5.\,1862 Wien – 21.\,10.\,1931 ebd.), \emph{Schriftsteller, Mediziner}!Vermächtnis. Schauspiel in drei Akten@\strich\emph{Das Vermächtnis. Schauspiel in drei Akten}|pw}}«, »\textsc{Freiwild\pwindex{Schnitzler, Arthur 15.\,5.\,1862 Wien – 21.\,10.\,1931 ebd.@\textsc{Schnitzler, Arthur} (15.\,5.\,1862 Wien – 21.\,10.\,1931 ebd.), \emph{Schriftsteller, Mediziner}!Freiwild. Schauspiel in 3 Akten@\strich\emph{Freiwild. Schauspiel in 3 Akten}|pw}}« {\kaufmannsund} »Garlan\pwindex{Schnitzler, Arthur 15.\,5.\,1862 Wien – 21.\,10.\,1931 ebd.@\textsc{Schnitzler, Arthur} (15.\,5.\,1862 Wien – 21.\,10.\,1931 ebd.), \emph{Schriftsteller, Mediziner}!Frau Bertha Garlan. Roman@\strich\emph{Frau Bertha Garlan. Roman}|pw}«
               zu finden, ſo erlaube ich mir darauf zu erwidern, daß mir dieſer Einwand nicht
               ganz ſtichhaltig erſcheint, weil dieſe Leute von untergeordneter Bedeutung ſind und
               ihr Beruf als ſolcher \uline{Nebenſache} iſt; hier aber ſteht
               gerade dieſer im Vordergrunde; überdies läßt gerade »Freiwild\pwindex{Schnitzler, Arthur 15.\,5.\,1862 Wien – 21.\,10.\,1931 ebd.@\textsc{Schnitzler, Arthur} (15.\,5.\,1862 Wien – 21.\,10.\,1931 ebd.), \emph{Schriftsteller, Mediziner}!Freiwild. Schauspiel in 3 Akten@\strich\emph{Freiwild. Schauspiel in 3 Akten}|pw}« wohl keinen Zweifel an Ihrer Geſinnung zu, trotzdem, wie Herr Schwarzkopf\pwindex{Schwarzkopf, Gustav 7.\,11.\,1853 Wien – 13.\,11.\,1939 ebd.@\textsc{Schwarzkopf, Gustav} (7.\,11.\,1853 Wien – 13.\,11.\,1939 ebd.), \emph{Schriftsteller}|pw} gelegentlich eines Gespräches zu
                Ihrer Verteidigung richtig {\pb}bemerkte, \introOben{}darin\introOben{}{ }\uline{nicht} der Doktor, ſondern der Offizier zum Helden
               wird. – Daß es Ihnen ganz ferne gelegen, in »Guſtel\pwindex{Schnitzler, Arthur 15.\,5.\,1862 Wien – 21.\,10.\,1931 ebd.@\textsc{Schnitzler, Arthur} (15.\,5.\,1862 Wien – 21.\,10.\,1931 ebd.), \emph{Schriftsteller, Mediziner}!Lieutenant Gustl. Novelle@\strich\emph{Lieutenant Gustl. Novelle}|pw}« den Offiziersſtand \uline{anzugreifen}, will
               ich gerne annehmen, ja, ich bin fest überzeugt, daß dies nicht der Fall geweſen iſt;
               Sie haben eben \uline{ohne} jede polemische Neben-Abſicht
               Ihre Geſinnung durchklingen lassen. \label{K_L04231-5v}\edtext{\textsc{\begin{otherlanguage}{french}Voilà tout\end{otherlanguage}}}{\lemma{\textnormal{\emph{Voilà tout}}}\Cendnote{\textnormal{französisch: das ist alles}}}\label{K_L04231-5}!
               Selbstverſtändlich hatten Sie das gute Recht dazu; und nicht nur hiezu,{ }ſondern auch
               zu einer Satire auf das Militär hätten Sie’s gehabt; ich wäre der letzte, dies
               bestreiten zu wollen, \uline{könnte} es nicht einmal, da
               ich ſelber Satiren im entgegengeſetzten Sinne geſchrieben habe. Anderſeits haben aber
               alle, die Anders denken, und mit dieſen ich, das Recht, Ihrer Oppoſition Ausdruck zu
               geben. Es handelt ſich nur darum, ob dis ſo geſchieht, daß \strikeout{I} die künſtleriſche Seite der Sache {\pb}hievon unberührt bleibe {\kaufmannsund} objektiv beurteilt werde. Daß dies bei der Neigung
               aller Leute zu ſubjektiver Beurteilung nur selten geſchehen dürfte, daß Sie von
               Anhängern meiner Geſinnung ebenſo heftig angegriffen als von ſolchen Ihrer Gesinnung
               geprieſen werden dürften, liegt eben in der Natur der Sache. Ich \uline{hoffe} darin eine Ausnahme zu machen, und da ich in Ermangelung \uline{feſter} Beziehungen zu irgend einem Blatte nicht ſicher
               wiſſen kann, ob ich auch Gelegenheit haben werde, dies schwarz auf weiß gedruckt zu
               bethätigen, ſo ſag’ ich’s eben hier, auf die Gefahr hin, daß es Ihnen – ſehr \uline{gegen} meinen Wunſch – \substVorne{}\textsuperscript{härter}\substDazwischen{}hart\substHinten{} klingt, da Sie ja in Ihrer Umgebung, unter Ihren Freunden{ }ſicherlich nur
                  Zuſti{\geminationm}ung finden werden, alſo nicht
               gewohnt ſind, ſozusagen direkt \strikeout{Ge} (dh. nicht durch
               Zeitungen \textsc{etc}) Gegenteiliges zu hören.\pend
           
\pstart
           Daß mein Urteil \strikeout{viell} beßer: meine Em{\pb}pfindung vielleicht anders wäre, wenn \strikeout{ich} mein Vater\pwindex{Sosnosky, ?? [Vater] @\textsc{Sosnosky, ?? [Vater]}, \emph{Offizier}|pwv}, meine beiden Großväter\pwindex{Sosnosky, ?? [Großvater väterlicherseits] @\textsc{Sosnosky, ?? [Großvater väterlicherseits]}, \emph{Offizier}|pwv}\pwindex{?? [Großvater mütterlicherseits von Theodor von Sosnosky] @\textsc{?? [Großvater mütterlicherseits von Theodor von Sosnosky]}|pwv} und
               mein Urgroßvater\pwindex{?? [Urgroßvater von Theodor von Sosnosky] @\textsc{?? [Urgroßvater von Theodor von Sosnosky]}|pwv}{ }\uline{nicht} Offiziere geweſen wären, das will ich nicht
               läugnen: es ka{\geminationn} aber kein Menſch aus ſeiner Haut
               heraus.\pend
           
\pstart
           Ihre Vermutung, es ſei das N. W. Tagbl.\pwindex{Neues Wiener Tagblatt@\emph{Neues Wiener Tagblatt}|pw} geweſen,
               das meine Beſprechung\pwindex{Sosnosky, Theodor von 4.\,1.\,1866 Budapest – 3.\,2.\,1943 Wien@\textsc{Sosnosky, Theodor von} (4.\,1.\,1866 Budapest – 3.\,2.\,1943 Wien), \emph{Schriftsteller}!Bücher und Zeitschriftenschau [Frau Bertha Garlan]@\strich\emph{Bücher und Zeitschriftenschau [Frau Bertha Garlan]}|pwv} »zu
               freundlich« gefunden hat, trifft \uline{nicht} zu; ich hatte
                  \strikeout{es}{ }ſie dieſem überhaupt nicht angeboten, ſondern
                  \label{T_L04231-1v}\edtext{einer}{\lemma{\textnormal{\emph{einer}}}\Cendnote{\textnormal{Er schreibt: »einerer«.}}}\label{T_L04231-1} ſehr bekannten
                  reichsdeutschen\oindex{Deutschland@\textbf{Deutschland}|pw}{ }Zeitſchrift\pwindex{?? [Zeitschrift, die Theodor von Sosnoskys Besprechung von »Bertha Garlan« zurückweist]@\emph{?? [Zeitschrift, die Theodor von Sosnoskys Besprechung von »Bertha Garlan« zurückweist]}|pwv}, deren Herauſgeber\pwindex{?? [Herausgeber einer reichsdeutschen Zeitschrift] @\textsc{?? [Herausgeber einer reichsdeutschen Zeitschrift]}|pwv} meinem Urteil
               über »Jung-Wien\orgindex{Jung Wien@Jung Wien|pw}« zuſtimmt, aber wie es ſcheint,
               nicht gleich mir, Ausnahmen gelten laßen will.\pend
           
\pstart
           Und nun,{ }ſehr geehrter Herr Doktor, ſchließe ich die lange Epiſtel und bitte Sie um
               Entſchuldigung, daß ich Sie ſo lange in Anſpruch geno{\geminationm}en; aber kürzer gieng’s halt nicht, wenn ich die Sachlage ſo klar machen ſollte,
               wie es mir in Ihrem {\kaufmannsund} in meinem Intereſse wünschenswert
               erſcheint. Hoffentlich nehmen Sie mir die Erörterung nicht übel und ſind nach wie vor
               meiner beſondren Wertſchätzung verſichert, mit der ich verbleibe Ihr ergebener\pend
           \pstart \spacefill\mbox{Theodor vSosnosky}\pend{}\selectlanguage{ngerman}\endnumbering\briefempfaengerindex{Schnitzler, Arthur@\textsc{Schnitzler, Arthur}!zzzSosnosky, Theodor von@\emph{von Theodor von Sosnosky}!1901-05-262@{26. 5. 1901}|)be}\mylabel{L04231h}
\begin{anhang}
\end{anhang}\newcommand{\dateiname}{L04231}\newcommand{\titel}{Theodor von Sosnosky an Arthur Schnitzler, 26. 5. 1901}\newcommand{\editorInnen}{Herausgegeben von Jahnke, SelmaMüller, Martin Anton}\input{../tex-inputs/latex-leseansicht-abspann}
